\chapter{数据治理、隐私保护与数字经济税收}

\begin{examguide}
\textbf{本章考情}：数据治理与数字税是 0258 论述题素材最密集的一章（\important{专硕·核心}）：三法框架与数据分类分级、双支柱方案的计算、隐私保护的制度比较均为高频命题。支柱二补足税与支柱一金额 A 的算例是计算题素材；隐私-效用权衡与避风港原则是简答题素材；隐私悖论与"被遗忘权"的表述辨析是辨析题常客。

\textbf{核心考点}：（1）隐私的经济学性质与隐私悖论（Acquisti 双命题）；（2）隐私-效用权衡的边际条件与最优规制水平、规制反对方四论点与 GDPR 悖论；（3）数据治理三法框架与数据分类分级；（4）知情同意的行为经济学（同意疲劳、选择过载）；（5）数据出境安全评估三路径；（6）算法治理与自动化决策规制；（7）GDPR 与中国个人信息保护制度的比较；（8）避风港原则、电商法连带责任与平台自我治理激励；（9）常设机构原则失效与 BEPS 转移定价；（10）数字服务税的税负转嫁；（11）双支柱方案（金额 A 分配、全球最低税补足税）与实施进展。

\textbf{本章主线}：隐私的经济学（12.1）→ 治理制度（12.2）→ 平台责任（12.3）→ 国际税收（12.4）。
\end{examguide}

第 11 章以成本结构为线索穿越了企业边界、市场边界与国家边界：数字化改变成本结构，成本结构改变组织与市场的边界。边界调整本身并不自动产生秩序：数据该不该跨境流动、隐私如何在效率与保护之间权衡、平台对生态内的行为承担何种责任、数字化利润在辖区之间如何分配，这些问题无法由市场自发解决，需要一套制度基础设施来回答。本章按"市场失灵、制度回应、执行责任、财政分配"的顺序展开：12.1 节用经济学语言回答"为什么要保护隐私"，12.2 节考察中国数据治理的法律体系与算法治理，12.3 节分析平台责任与生态治理的激励结构，12.4 节处理数字经济对国际税收秩序的冲击与双支柱方案。

\section{隐私经济学}

治理制度的设计要先回答一个前置问题：被治理的对象在经济上是什么性质？隐私治理的对象是个人信息，其经济学性质决定治理工具的选择。本节分三步：先刻画隐私作为商品的特殊性质，再推导消费者以隐私换服务的最优条件，最后把外部性纳入规制水平的分析。

\subsection{隐私的经济学性质}

隐私的经济学分析起点，是将隐私视为一种\important{可交易的商品}：消费者以隐私信息换取免费或低价服务，这是第 5 章三方市场的另一面，也是免费商业模式得以运转的对价。但隐私不是普通商品，它有三个特殊性质，使这桩交易系统性地偏离标准商品交易的逻辑。其一，\textbf{信息不对称}：数据加工的链条长且不透明，一条浏览记录经过收集、整合、推断、出售，最终变成什么用途、产生什么价值，消费者几乎无法追踪；其二，\textbf{外部性}：个体的数据披露可能暴露他人的信息，社交网络数据中的朋友信息、群体画像对相似个体的推断都是例子，个体决策不会内化对他人的隐私损害；其三，\textbf{不可逆性}：数据一旦披露，复制与传播不可追溯，事后撤回的效力有限，这一性质与第 7 章讨论的数据要素属性同源。

这三个性质叠加的结果，是隐私交易在真实市场中的表现与教科书上的商品交换大不相同，最著名的经验事实就是\important{隐私悖论}。理解这一悖论是理解整个隐私治理制度的入口：如果消费者的行为系统性地偏离自己的态度，那么"尊重消费者选择"就不再是充分的政策原则，制度需要处理的不是选择本身，而是选择背后的决策环境。

\begin{definition}[隐私悖论（Privacy Paradox）]
隐私悖论是指消费者在态度调查中声称高度重视隐私保护，却在现实选择中以微小对价主动放弃隐私信息，态度与行为系统性背离的现象。
\end{definition}

悖论最早来自观察性证据，随后被实验反复证实。在 Acquisti 团队的经典实验中，声称注重隐私的受访者超九成以一张披萨代金券为对价披露了朋友的邮箱地址（Acquisti 等，2009），隐私声明的强度与实际披露行为几乎不相关。悖论的解释不是"消费者的偏好自相矛盾"：同一个人的确既看重隐私、又想占便宜，真正的问题是决策环境使这两个偏好无法在同一个决策中同时兑现。Acquisti 把这种系统性偏离归结为两个命题。\textbf{信息不对称}：消费者对披露的真实代价缺乏认知，不知道自己交出去的数据值多少钱、会带来什么后果；\textbf{跨期不一致性}：隐私损失在远期、服务收益在当下，消费者的贴现结构会系统性地压低远期损失的权重，使他在决策时低估隐私的代价。第一个命题靠制度强制披露来缓解（告知义务），第二个命题则要先形式化，才能看清政策的着力点。

\begin{derivation}{跨期不一致性与过度披露}{}
设消费者在第 0 期决定隐私披露量 $d$，披露带来即期收益 $u(d)$（免费服务、个性化便利），满足 $u'>0$、$u''<0$，即收益递增但边际递减；隐私损失在下一期兑现，损失函数为 $c(d)$，满足 $c'>0$、$c''>0$，即损失递增且边际加速。消费者的贴现结构采用拟双曲贴现 $(\beta,\delta)$：$\delta\in(0,1)$ 是常规贴现因子，负责把未来效用逐年折现；$\beta\in(0,1)$ 是当下偏差参数，刻画"现在与未来之间"的额外权重折扣。同一单位效用，只要它落在未来，消费者对它的重视程度就要先乘 $\beta$ 再按 $\delta$ 折现，这是行为经济学对自控问题最常用的刻画方式。
第一步，决策期的最优披露。第 0 期做决策时，服务收益立即兑现，按权重 $1$ 计入；隐私损失要到第 1 期才发生，先被 $\beta$ 打折、再被 $\delta$ 折现，因此目标函数是
\[
u(d) - \beta\delta\, c(d).
\]
对 $d$ 求导并令一阶导数为零：
\[
u'(d^*) = \beta\delta\, c'(d^*).
\]
等式的含义：决策者会把披露量一直推到"最后一单位披露的即期边际收益恰好等于被双重打折的边际损失"为止；边际损失的权重 $\beta\delta$ 越低，他愿意接受的披露量就越高。
第二步，无扭曲的事后基准。作为比较基准，设想一个收益与损失都不被折现的参照系，即事后充分权衡、不受 $\beta$ 与 $\delta$ 折损的选择，其最优披露 $d^{FB}$ 满足
\[
u'(d^{FB}) = c'(d^{FB}),
\]
即边际收益等于不打折的边际损失。$d^{FB}$ 就是消费者事后回顾时真正想要的披露水平，也是问卷中"重视隐私"的态度所对应的选择。
第三步，过度披露的比较。为了确定两个最优量的大小关系，构造辅助函数 $F(d) = u'(d) - \beta\delta c'(d)$。由 $u''<0$、$c''>0$ 可知 $F$ 是 $d$ 的严格减函数，因此 $F$ 的零点唯一，即 $d^*$ 唯一。把事后基准 $d^{FB}$ 代入，并利用第二步的一阶条件 $u'(d^{FB}) = c'(d^{FB})$：
\[
F(d^{FB}) = u'(d^{FB}) - \beta\delta c'(d^{FB}) = (1-\beta\delta)\, c'(d^{FB}) > 0.
\]
上式说明：在事后基准的位置上，决策期的"净边际效用"仍然是正的，决策者仍想继续披露，只有越过 $d^{FB}$ 继续向右，$F$ 才会降到零。因此
\[
d^* > d^{FB}:
\]
\important{当下偏差使决策期的披露量系统性超过事后最优水平}。
\end{derivation}

把两步最优放在一起看，隐私悖论的谜底就解开了。消费者在问卷里表达的"重视隐私"，对应的是无扭曲的事后基准 $d^{FB}$：一个收益与损失同等权重、充分权衡过的选择；而在真实选择中轻易放弃隐私，对应的是决策期最优 $d^*$：一个对当期收益偏好倾斜的选择。两者并不矛盾，它们是同一个人在不同时间点上的两个最优选择，悖论的根源不是偏好混乱，而是贴现结构的时间不一致性。沿着这条线索还能看清一个常被忽略的事实：即便是指数贴现的消费者（$\beta=1$，不存在当下偏差），决策时也会因为未来损失被 $\delta$ 逐年折现而比事后基准披露得更多，这属于跨期选择的一般性质，来自理性的时间偏好；拟双曲贴现的贡献是在这个基础上再加一层 $\beta$ 的额外折扣，使过度披露进一步放大，并且制造出"决策时觉得划算、事后后悔"的时间不一致体验。

这一机制直接给出政策的着力点。既然过度披露的根源在贴现结构，仅靠"告知风险"就不够：告知改变的是消费者对 $c(d)$ 的认知，针对的是信息不对称命题；而跨期不一致性命题要求制度替"事后自我"守住底线。默认规则（默认不收集、最小化收集）直接降低决策期的披露起点，消费者不必在每个场景下重新战胜自己的当下偏差；删除权与撤回同意权则给"事后自我"提供了纠错通道，后悔的消费者可以把数据拿回来。这两类工具正是《通用数据保护条例》（GDPR）与《个人信息保护法》（以下简称个保法）的两大支柱：告知同意机制针对信息不对称，数据权利与默认规则针对跨期不一致性，12.2 节展开。

\subsection{隐私换服务的边际条件}

免费商业模式（第 5 章）的本质，是消费者以数据"支付"服务：不付现金，但要交出个人信息，平台再用这些数据做定向广告与产品改进。第 10 章讨论数据征信时已经触及这个权衡的轮廓：评估精度的边际收益以隐私的边际成本为对价（第 10 章）。本节把这个权衡完整地写成模型。要回答的问题是：消费者面对"多披露一点数据、多获得一点服务"的替代关系，最优的披露边界在哪里？政策又如何移动这条边界？

\begin{derivation}{隐私换服务的边际条件}{}
设消费者使用数字服务的效用为 $u(s,d)$，其中 $s$ 为服务价值，$d$ 为披露的隐私数据量，满足 $\partial u/\partial s > 0$、$\partial u/\partial d < 0$：服务越好效用越高，而披露本身直接降低效用，因为隐私暴露带来风险成本与心理成本。服务提供者以数据为生产投入（定向广告、产品改进），服务价值随数据披露上升：$s = s(d)$，$s'>0$ 且 $s''<0$，后一个条件表示数据对服务的边际贡献递减，即多交一份数据的服务增量越来越小。
第一步，一阶条件。披露量 $d$ 通过两条渠道影响效用：一条正渠道，数据让服务升级；一条负渠道，披露直接损害隐私。总效用对 $d$ 的全导数是两项之和，最优披露要求全导数为零：
\[
\frac{du}{dd} = \frac{\partial u}{\partial s} s'(d) + \frac{\partial u}{\partial d} = 0
\;\Rightarrow\;
\underbrace{\frac{\partial u}{\partial s} s'(d^*)}_{\text{服务的边际收益}} = \underbrace{-\frac{\partial u}{\partial d}}_{\text{隐私的边际成本}}.
\]
等式的含义：最优披露量停在服务的边际收益恰好等于隐私的边际成本的位置。左边是"多披露一点数据带来的服务改善"，右边是"多披露一点数据造成的隐私损失"，两者在边际上打平，再披露任何一单位都会得不偿失。
第二步，二阶条件。二阶条件保证上述解是极大值点而非极小值点：服务的边际收益随 $d$ 递减（$s''<0$），隐私的边际成本随 $d$ 递增（披露越多，追加披露的损失越敏感），两条曲线的斜率方向相反，交点唯一。
第三步，数值算例。算例把模型具体化。取效用函数 $u(s,d) = 4\sqrt{s} - d$：服务带来的正效用递增但边际递减，披露带来的负效用随披露量线性上升；服务生产函数 $s(d) = \sqrt{d}$ 满足边际贡献递减。代入得总效用 $U(d) = 4d^{1/4} - d$，一阶条件为
\[
U'(d) = d^{-3/4} - 1 = 0 \;\Rightarrow\; d^* = 1.
\]
在 $d^* = 1$ 处，服务的边际收益 $d^{-3/4} = 1$，恰好等于隐私的边际成本 $1$，与一般模型的结论一致。
第四步，政策比较静态。政策实验：假设法律赋予消费者删除权，平台被要求删除数据后，消费者披露数据的预期成本上升，等价于隐私的边际成本由 $1$ 升至 $2$。新的最优披露满足 $d^{-3/4} = 2$，解得
\[
d^{**} = 2^{-4/3} \approx 0.40,
\]
披露量较 $d^* = 1$ 下降约六成。\important{数据权利的实质是抬高隐私的边际成本，使均衡披露量左移}。
\end{derivation}

均衡条件的含义直观：披露过少，服务的边际收益没有被充分利用，消费者白白损失了本可以免费获得的服务；披露过多，隐私的边际成本超过收益，每一单位追加的披露都在净亏损。\mn{边际条件白话：多交一份数据的服务好处等于多交一份数据的隐私代价，两边打平就是最优；披露少了亏服务，披露多了亏隐私。}最优披露量位于两条边际曲线的交点，任何一边的偏离都意味着消费者的效用没有最大化。

政策的两个切入点由此明确。第一个切入点是\textbf{提高边际成本}：赋予消费者数据权利（知情权、删除权、撤回同意权）使隐私披露的预期成本上升，算例显示边际成本翻倍使披露量下降约六成，权利的力度可以直接换算成披露量的变化。第二个切入点是\textbf{降低信息不对称}：强制告知与隐私政策的简明化要求使消费者对 $s'(d)$ 与隐私风险形成真实认知，否则边际收益与边际成本的比较根本无从发生，消费者在不知道函数形状的情况下不可能解出最优披露量。两个切入点分别对应 GDPR 与个保法的两大支柱：数据权利体系与告知同意机制。这一框架同样是第 10 章数据征信权衡的一般化：征信精度的边际收益必须以隐私的边际成本为对价，任何以数据为投入的业务（征信、广告、推荐）都可以代入本节的边际条件，区别只在于效用函数的具体形状。

\subsection{最优规制水平}

上一节的最优化是消费者的个人决策，消费者自己承担披露的全部成本与收益。但隐私披露的损害并不全部由披露者自己承担：朋友信息被连带披露、群体画像波及相似个体，隐私具有负外部性。一旦私人决策的成本收益与社会的不一致，市场自发的结果就不再是社会想要的，规制的必要性与强度边界就成为一个可计算的问题。

\begin{derivation}{隐私规制的外部性框架}{}
设企业数据采集量为 $d$，私人边际收益为 $MB(d)$（数据驱动的收入，如定向广告收入随数据量增加而上升），满足 $MB'<0$，即数据的边际收入递减；私人边际成本为 $MC(d)$（合规成本、声誉风险），满足 $MC'>0$。企业决策时不承担的隐私损害为 $E(d)$，满足 $E'>0$：企业每多采集一单位数据，消费者整体额外承受一份隐私损失，而企业不为这份损失付费。
第一步，私人最优。企业最大化自己的净收益，边际上要求私人边际收益等于私人边际成本：
\[
MB(d_p^*) = MC(d_p^*).
\]
第二步，社会最优。社会计划者把外部成本计入账本：社会希望边际收益等于私人边际成本与外部边际成本之和：
\[
MB(d_s^*) = MC(d_s^*) + E'(d_s^*).
\]
第三步，偏离方向。在私人最优处检验社会最优条件是否成立：
\[
MB(d_p^*) = MC(d_p^*) < MC(d_p^*) + E'(d_p^*).
\]
私人最优的边际收益只够覆盖企业的私人成本，覆盖不了社会成本，社会视角下"最后一单位采集"是亏损的。由于 $MB$ 随 $d$ 递减、$MC + E'$ 随 $d$ 递增，要让边际收益提高到与社会成本打平，采集量必须往回收：
\[
d_s^* < d_p^*,
\]
\important{私人最优的采集量超过社会最优}，隐私保护规制（数据最小化原则、目的限制原则）的经济学定位即对这一负外部性的矫正。
\end{derivation}

\begin{figure}[htbp]
\centering
\begin{tikzpicture}[>=Stealth, x=1.3cm, y=1.3cm]
\draw[->, very thick] (0,0) -- (0,7.3);
\draw[->, very thick] (0,0) -- (7.6,0);
\node[font=\small, rotate=90] at (-0.7,3.5) {边际收益与边际成本};
\node[font=\small, below=6pt] at (3.8,0) {数据采集量 $d$};
\draw[main, very thick] (0,7) -- (7,0) node[pos=0.28, above=3pt, sloped, font=\small, main] {边际收益 $MB$};
\draw[second, very thick] (0,1) -- (7,2.4) node[pos=0.35, below=3pt, sloped, font=\small, second] {私人边际成本 $MC$};
\draw[winered, very thick] (0,2.2) -- (7,3.6) node[pos=0.28, above=3pt, sloped, font=\small, winered] {社会边际成本 $MC+E'(d)$};
\draw[dashed, structurecolor, thick] (4,0) -- (4,7.3);
\draw[dashed, structurecolor, thick] (5,0) -- (5,7.3);
\fill[main] (4,3) circle (1.5pt);
\fill[main] (5,2) circle (1.5pt);
\node[font=\small, structurecolor, fill=white, inner sep=2pt, above=8pt] at (4,3) {社会最优 $d_s^*$};
\node[font=\small, structurecolor, fill=white, inner sep=2pt, below=8pt] at (5,2) {私人最优 $d_p^*$};
\end{tikzpicture}
\caption{隐私规制的私人最优与社会最优}
\label{fig:privacy-ext}
\end{figure}

图~\ref{fig:privacy-ext} 给出了两种最优的比较：社会边际成本曲线位于私人边际成本曲线上方，二者的垂直距离正是外部成本 $E'(d)$，社会最优采集量因此位于私人最优的左侧。规制强度的边界同样由社会最优条件给出：过度规制把采集量压到 $d_s^*$ 以下，数据驱动的创新受损，规制同样有代价。规制的目标因此不是消除采集，而是使企业内化 $E'(d)$，让企业的私人最优自动对齐社会最优。

内化的工具有三类。罚款与公益诉讼把外部成本直接写回企业的成本函数，企业每多采集一单位数据的预期罚金就是它对社会的赔偿；集体赔偿制度降低单个消费者的维权成本，解决外部性案件中"受害者众多、单笔损失微小"导致的搭便车问题；行业标准（最小化收集、默认关闭、目的限制）则把内化前置到产品设计环节，从源头压缩过度采集的空间。实证方面，对欧盟的评估研究显示 GDPR 生效后第三方追踪与广告请求显著减少，说明规制确实压缩了采集量；但对企业融资与网站流量的影响方向在文献中存在分歧，提示规制强度存在真实的经济代价，社会最优条件的两端都需要证据支撑。中国的做法是行政通报：工业和信息化部自 2021 年起分批通报侵害用户权益的 App，问题集中于违规收集个人信息与强制索权（工业和信息化部，2021--2024），行政通报的成本远低于罚款程序，同样把声誉损失写进企业的成本函数。

\textbf{反对规制的一方}同样有系统的论证，论述题"你是否支持隐私规制"要求两面都答。其论证有四。其一，\textbf{信号传递受阻}：芝加哥学派认为隐私保护削弱了"好消费者"向市场发送可信信号的机制，自愿披露本可替代外部评估，规制提高了甄别成本；其二，\textbf{商业模式创新受阻}：广告支持的免费互联网依赖数据变现，规制压缩数据供给，免费模式的资金来源被切断，消费者以付费形式承担损失；其三，\textbf{知情同意成为进入壁垒}：大平台拥有法务与合规团队，能把合规成本摊薄，初创企业难以承担逐项同意的合规负担，规制事实上保护了在位者（第 6 章的市场势力逻辑在制度层面重现）；其四，\textbf{GDPR 的"干净数据"悖论}：Aridor 等（2020）指出，知情同意挤掉了不完美的信息隐藏，留存的数据更"干净"、更准确，反而使企业更易追踪与画像，规制对隐私的保护作用被数据质量的提升部分抵消。\mn{隐私规制的反方论点速记：信号、模式、壁垒、悖论。正论用外部性框架（$d_s^* < d_p^*$），反论报四论点，收尾"规制强度由 $E'(d)$ 的可度量性决定"，两面都要会写。}一个更尖锐的理论结果来自 Acquisti 等（2016）：即使产权配置完美清晰，垄断厂商也可以凭借"微小报价"（一张小额优惠券）套出消费者保留价格，\important{清晰的产权并不能自动保护隐私，市场势力才是隐私威胁的独立来源}，这为反垄断与隐私保护的交汇（第 6 章）提供了理论基础。正反双方的共同落点是：\important{隐私规制的强度边界由外部成本 $E'(d)$ 的可度量性决定}，证据充分处用外部性框架矫正，证据不足处的过度规制同样有代价。

\section{数据治理法律体系}

隐私经济学给出的是"为什么需要治理"的一般原理，治理的落地需要制度供给。中国的数据治理以三部法律为骨架、以配套规章为血肉，总体思路可以概括为：分对象立法、分风险配置义务、分场景设计规则。本节沿这一思路展开：12.2.1 节给出三法框架与数据分类分级，12.2.2 节分析知情同意的经济学，12.2.3 节考察数据出境的阶梯化安排，12.2.4 节把治理延伸到算法层，12.2.5 节与 GDPR 做制度比较。

\subsection{三法框架与数据分类分级}

治理制度的第一步是确定"管谁、管什么"。中国数据治理的三部法律以对象为界分工，形成"网络基础设施、数据本体、个人信息"三层递进的结构。《网络安全法》（2017 年 6 月施行）管\textbf{网络运行安全}：关键信息基础设施、等级保护、网络运营者义务；《数据安全法》（2021 年 9 月施行）管\textbf{数据本体}：数据分类分级、数据安全审查、跨境流动安全评估；《个人信息保护法》（2021 年 11 月施行）管\textbf{个人数据}：个人信息处理规则、个体权利、跨境提供规则。\mn{三法分工记忆：网安法管"网络"，数安法管"数据"，个保法管"人"。三者关系：网安法是基础设施前提，数安法是数据治理总框架，个保法是个人数据的专门法。}三者的关系是递进的：网安法提供基础设施前提，数安法建立数据治理总框架，个保法是个人数据领域的专门法。图~\ref{fig:3laws} 给出了这一分工结构。

\begin{figure}[htbp]
\centering
\begin{tikzpicture}[>=Stealth,
    law/.style={rectangle, rounded corners=2mm, draw=main, thick, fill=main!8, align=center, font=\small, minimum width=34mm, minimum height=11mm},
    obj/.style={rectangle, rounded corners=2mm, draw=second, thick, fill=second!8, align=center, font=\small, minimum width=28mm, minimum height=11mm},
    arr/.style={->, very thick},
    note/.style={font=\small, color=structurecolor}]
\node[law] (L1) at (0,0) {《网络安全法》\\2017 年施行};
\node[obj] (O1) at (5.6,0) {网络运行安全};
\draw[arr] (L1) -- node[above=2pt, font=\small] {管} (O1);
\node[note, right=5mm of O1] {基础设施前提};
\node[law] (L2) at (0,-2.6) {《数据安全法》\\2021 年施行};
\node[obj] (O2) at (5.6,-2.6) {数据本体};
\draw[arr] (L2) -- node[above=2pt, font=\small] {管} (O2);
\node[note, right=5mm of O2] {数据治理总框架};
\node[law] (L3) at (0,-5.2) {《个人信息保护法》\\2021 年施行};
\node[obj] (O3) at (5.6,-5.2) {个人信息};
\draw[arr] (L3) -- node[above=2pt, font=\small] {管} (O3);
\node[note, right=5mm of O3] {个人数据专门法};
\end{tikzpicture}
\caption{数据治理三法的分工}
\label{fig:3laws}
\end{figure}

三法的共同逻辑是\textbf{风险分级治理}：按数据的重要程度与敏感程度配置差异化的保护义务。《数据安全法》第二十一条确立数据分类分级保护制度，将数据分为一般数据、重要数据与核心数据，核心数据关系国家安全、国民经济命脉、重要民生与重大公共利益，实施最严格保护；《个人信息保护法》第二十八条区分一般个人信息与敏感个人信息，生物识别、宗教信仰、特定身份、医疗健康、金融账户、行踪轨迹等信息属于敏感个人信息，处理敏感个人信息须取得单独同意。\mn{易混点：数据分类分级（一般、重要、核心数据）出自《数据安全法》，管数据多重要；敏感个人信息（生物识别、金融账户等）出自《个人信息保护法》，管人有多敏感。场景题先定位是哪部法。}分级治理的经济学含义：安全投入是稀缺资源，按风险配置保护强度（高风险高保护）避免"一刀切"的过度合规成本；反过来，对低风险数据的过度保护同样有代价，中小企业的合规负担被无谓抬高，数据流通的效率受损。\important{分级分类是数据治理的成本有效性设计}。这与第 7 章数据二十条的产权分置构成配套：要素侧"三权分置"促流通，治理侧"分级分类"保安全，两头并进是数据要素市场的中国式设计。

\subsection{知情同意的经济学}

三部法律确立了权利框架，但制度效力取决于信息处理的合法性基础能否真实运行。个人信息处理的合法性基础是\important{告知—同意}：处理者告知处理规则，个人同意后处理才可进行。这个机制的设计意图是让消费者自己把关：知情之后的选择，理论上接近 12.1.2 节的最优披露决策。但行为经济学提出了一个尖锐的质疑：如果充分知情的成本超过了消费者的承受能力，"同意"还能承载真实的意思表示吗？本节用注意力预算模型把"同意疲劳"形式化，回答这个质疑。

\begin{derivation}{知情同意的注意力预算模型}{}
设消费者面对 $N$ 份数字服务的隐私政策，第 $i$ 份政策的完整阅读成本为 $c_i > 0$（时间与认知成本：读懂一份隐私政策需要逐条理解法律术语），阅读带来的知情收益为 $b_i$（知情后可以避免的预期隐私损失：避开过度收集的服务、放弃不利条款）。
第一步，单份政策的阅读决策。消费者阅读第 $i$ 份政策，当且仅当阅读的净收益非负：
\[
b_i \ge c_i.
\]
如果读懂一份政策要花的时间成本高于它能为消费者避免的损失，理性的选择是不读。
第二步，注意力预算约束。阅读决策不能只看单份政策：消费者同时面对 $N$ 份政策，注意力总量有限。完整知情的可行条件是所有政策的阅读成本之和不超过注意力预算 $\bar{C}$：
\[
\sum_{i=1}^{N} c_i \le \bar{C}.
\]
当条款数量与冗长程度使 $\sum_i c_i > \bar{C}$ 时，完整阅读在物理上不可行，"全部知情"从一开始就不是一个选项。
第三步，同意疲劳。即使注意力充足，消费者也没有动机逐份阅读：标准化条款重复出现，读第二份政策的增量收益远低于第一份（$b_i$ 随 $i$ 递减），而阅读成本近似不变，存在阈值 $i^* = \max\{i : b_i \ge c_i\}$；对 $i > i^*$ 的政策，阅读的净收益为负，理性选择是"不看而同意"。这就是同意疲劳：不是消费者懒，而是在边际上继续阅读不划算。
第四步，退化均衡。企业一侧同样在优化：冗长的条款降低了消费者知情的概率，同时把企业的告知义务形式上履行完毕，因此企业有动机提高条款密度与条款数量（$c_i$ 上升、$N$ 上升），把阈值 $i^*$ 压向零。结果是知情同意的实际比例 $q = i^*/N$ 趋近于零：\important{当完整知情的成本超过注意力预算，同意退化为习惯性点击，失去筛选功能}。市场自发的结果是：设计良好的同意机制在成本结构上就注定失灵。
\end{derivation}

模型的结论与行为经济学的证据一致，三个效应共同解释了"同意失灵"。\textbf{有限注意力}：消费者的认知资源是稀缺的，面对冗长条款的最优反应是跳过阅读，注意力在经济上是需要配置的资源而非无限供给；\textbf{选择过载}：同时呈现的选项与政策数量越多，消费者越依赖默认选项与启发式决策，决策质量反而下降；\mn{同意疲劳论述题素材：条款冗长使"同意"沦为形式；个保法回应三件套：单独同意（敏感信息）、撤回同意权、自动化决策拒绝权；改进方向：降知情成本、提默认保护，而非加长文本。}\textbf{零价格效应}：免费服务的"零价格"标签使消费者低估隐私对价，"没花钱"的错觉掩盖了"用数据付费"的事实，这正是第 5 章注意力经济的微观基础。

模型同时指明制度改进的方向，并且说明什么方向是错的。错误方向是继续加强"同意"的形式：更长的政策文本只会进一步压低 $i^*$，使更多消费者放弃阅读。正确方向有两个。\textbf{降低知情成本}：简明告知、图标化提示、隐私政策的标准化摘要，把 $c_i$ 降下来，让阅读重新变得划算；\textbf{提高默认保护}：默认不收集、最小化收集、目的限制，把不阅读的消费者也保护起来，使"同意疲劳"不再等同于"裸奔"。更根本的调整是改换合法性基础：GDPR 与个保法都在"同意"之外保留了合同必要、法定义务等合法性事由，避免一切处理都以一份失灵的同意为凭据。个保法对敏感个人信息要求单独同意、赋予撤回同意权与自动化决策拒绝权，其意图正是绕开"一揽子同意"的退化均衡。

\subsection{数据出境安全评估}

第 11 章在数字贸易规则的讨论中留下一个权衡：数据自由流动的效率诉求与数据本地化的安全诉求之间存在张力（第 11 章）。中国的制度选择是让出境成本随风险阶梯化，用成本差异引导数据分级流动。数据跨境流动的制度安排呈现\important{三路径}结构：\textbf{安全评估}：重要数据与关键信息基础设施运营者的个人信息出境，以及处理个人信息达一定规模的情形（如处理 100 万人以上个人信息，或累计向境外提供 10 万人以上个人信息、1 万人以上敏感个人信息），须经网信部门安全评估（国家互联网信息办公室，2022）；\textbf{标准合同}：与境外接收方签订标准合同并备案；\textbf{认证}：经专业机构进行个人信息保护认证。2024 年施行的《促进和规范数据跨境流动规定》对不涉及重要数据的一般情形适度简化程序，安全评估聚焦高风险场景。\mn{出境三路径速记：安全评估（如处理 100 万人以上个人信息，或累计出境 10 万人以上个人信息、1 万人以上敏感个人信息，国家互联网信息办公室 2022 年口径）、标准合同、认证；风险越高，路径越严。}这一校准的经济逻辑与分级治理一致：一般数据出境的企业合规成本如果过高，跨境电商、技术服务外包等日常商业活动会被无谓拖累，简化的对象正是低风险情形。

三路径的经济学含义是交易成本的阶梯化：数据的重要性越高、风险越大，出境的合规成本越高，以成本差异引导数据分级流动，\important{出境制度的实质是给数据流动定价，价格随风险上升}，在自由流动的效率与安全可控的要求之间建立分层折中。这一安排同时是第 11 章数字贸易规则的国内制度基础：对外的规则谈判（DEPA、CPTPP 中的数据流动条款）与对内的出境安排共享同一套风险逻辑，对外承诺开放的部分对应境内低风险路径，对外保留安全例外的部分对应境内评估路径。

\subsection{算法治理}

数据治理的第三层是算法层：数据的风险经过算法转化为决策，规制的对象因此从"数据本身"延伸到"数据如何被使用"。2022 年施行的算法推荐管理规定把算法纳入制度框架，其核心安排与第 5 章的价格歧视理论直接呼应。《互联网信息服务算法推荐管理规定》（2022 年 3 月施行）确立三项核心制度。\textbf{算法备案}：具有舆论属性或社会动员能力的算法推荐服务提供者应当在提供服务之日起 10 个工作日内履行备案手续，公示算法的基本原理、目的意图与运行机制，并实行分类分级管理；\textbf{不得差异化定价}：第二十一条禁止利用算法在交易价格等交易条件上实施不合理的差别待遇，这是大数据杀熟的直接规制依据。与第 5 章的衔接在于行为价格歧视的实现条件：数据提供画像素材，算法实施定价决策，市场势力保证消费者无法逃离，三者缺一不可。数据治理解决了数据环节的合规问题，反垄断执法解决市场势力环节（第 6 章），算法治理补上的正是"算法"环节：切断算法对不合理差别的实施能力；\textbf{用户选择权}：提供不针对个人特征的选项或便捷的关闭方式，使拒绝个性化成为可行选项。个保法第二十四条进一步规定：自动化决策应当保证透明度和结果公平，不得实行不合理的差别待遇，个人有权拒绝仅通过自动化决策作出的决定。\mn{算法治理三件套：备案（管进入）、透明（管过程：公示机制、解释义务）、选择（管结果：个性化关闭、自动化决策拒绝）。}生成式人工智能的出现把算法治理推向纵深：《生成式人工智能服务管理暂行办法》（2023 年 8 月施行）对具有舆论属性或社会动员能力的大模型服务实施安全评估与算法备案双重义务。

算法治理的经济学定位由此清晰：算法是市场势力的放大器，同样的数据在不同算法下产生不同的定价与推荐后果，算法本身的复杂性又使外部监督难以穿透，因此规制对象是决策行为而非技术本身，\important{算法治理的经济学实质是把算法的决策后果纳入责任体系}。这一框架与第 6 章算法合谋的讨论同源：合谋、歧视、操纵都是算法决策的外部后果，事前备案让监管者看见算法，事后责任让平台承担后果，两者共同构成算法问责的制度闭环。

\subsection{GDPR 与中国制度比较}

中国的数据治理并非孤立存在。全球数据治理的两大范式是欧盟的权利本位与中国的发展安全平衡，理解二者的差异与共性，是把握全球数据治理格局的起点。《通用数据保护条例》（GDPR）于 2018 年 5 月生效，三大制度特征构成其标杆地位。\textbf{域外管辖}：第三条对向欧盟境内主体提供商品或服务、或监控其行为的处理者行使管辖，无论企业位于何地，即所谓"长臂管辖"；\textbf{被遗忘权}：第十七条赋予数据主体要求删除其个人数据的权利，处理者还有义务通知其他处理者删除链接与副本；\textbf{威慑型罚款}：最高 2000 万欧元或全球年营业额的 4\%，取较高者。\mn{辨析题常客："中国《个人信息保护法》确立了被遗忘权"是错误表述，个保法的对应制度是删除权，且附有法律、行政法规规定的保存期限等限制；"GDPR 的被遗忘权无任何限制"同样是错误表述，言论自由等豁免情形构成其边界。}个保法在结构上与之同构：第三条设置域外适用（以向境内自然人提供产品或服务为目的、或分析评估境内自然人行为的处理活动适用本法）；个体权利体系涵盖知情权、决定权、查阅复制权、更正权、删除权、撤回同意权；罚款上限为 5000 万元或上一年度营业额的 5\%。

比较的学理要点在范式差异。欧盟将隐私定位为基本权利，源于欧洲战后的人格尊严传统，GDPR 是这一权利传统的法典化：保护水平由权利清单统一规定，不因场景、行业、数据量而变，制度重心在个体的对抗性权利，以高额罚款维持威慑。代价是合规成本高，数据流通与创新的空间被压缩，中小企业的负担尤其突出。中国将数据治理嵌入发展与安全的双重目标，个保法与数据要素市场建设（第 7 章数据二十条）互为表里：既承认个体权利，又要为数据流通保留制度通道，制度重心在风险分级与流通促进的平衡。代价是层级之间的边界需要精细界定，否则企业面临合规不确定性。两者共享"域外管辖 + 高额罚款"的威慑结构，且都以数据权利体系与告知同意机制为支柱（12.1.2 节）。\important{两大范式的差异不在是否保护，而在保护的实现方式：欧盟以权利清单为核心，中国以分级治理为核心}。这就是制度比较的基本结论。

\section{平台责任与数字生态治理}

制度文本的执行落在平台。平台是数字生态事实上的治理者：它制定规则、分配流量、裁决纠纷，治理的成败取决于平台的激励是否与社会目标对齐。本节先考察平台责任的制度边界，再分析平台自我治理的激励机制。

\subsection{避风港原则与平台责任}

平台治理的第一个问题是责任边界：平台对生态内的侵权内容承担何种责任？责任过重，事前审查义务扼杀内容供给；责任过轻，侵权泛滥。制度设计的核心是审查成本的分配，\textbf{避风港原则}正是这一成本分配的制度化表达。

\begin{definition}[避风港原则（Safe Harbor）]
避风港原则是指平台对用户发布的内容不承担事前审查义务，但在收到权利人通知后须及时移除侵权内容（"通知—删除"规则），否则就损害扩大的部分承担责任的制度安排。
\end{definition}

避风港原则的经济学逻辑是成本分配：平台的事前审查成本极高，海量内容的逐一审查不可行，强制事前审查等于关闭内容供给；将审查义务后置为"通知后删除"，把治理成本分配给最有条件承担的双方：权利人最了解自己的作品被谁侵权（信息优势），平台最有技术能力执行删除（执行优势），两者的合作以最低成本实现侵权处理。避风港的边界是\textbf{红旗原则（Red Flag）}：当侵权事实像红旗一样明显时（知名作品被整段上传），平台不能以未收到通知抗辩，明知或应知即不免责。红旗原则防止平台"选择性失明"：把眼睛蒙上就能躲进避风港的话，避风港就成了纵容侵权的制度漏洞。\mn{避风港三件套：通知—删除（不知情则免责）、红旗原则（应当知情不免责）、电商法连带责任（知道或应知未处置则连带）。}中国法的对应安排：《信息网络传播权保护条例》（2006）确立"通知—删除"规则；《电子商务法》（2019 年施行）第三十八条将责任延伸至消费者权益：平台知道或应当知道平台内经营者侵害消费者合法权益，未采取必要措施的，与经营者承担连带责任；对关系消费者生命健康的商品或服务，未尽审核义务或安全保障义务的，承担相应责任。

数字时代对避风港原则的挑战来自算法分发。避风港的成立前提是平台的中立通道地位：内容由用户上传，平台只是载体。当平台以算法决定内容的可见性（推荐而非被动存储），平台的职能就越过了通道的边界：它决定用户看到什么、不看到什么，这与编辑决定报纸版面没有本质区别，其角色更接近出版者而非书店。角色变了，责任标准随之移动。《互联网信息服务算法推荐管理规定》（2022 年施行）要求算法推荐服务提供者对违法信息承担管理责任：\important{避风港原则在算法分发场景收缩，平台从"通道"走向"编辑"}。

\subsection{平台自我治理的激励}

责任边界之外，更深刻的问题是激励：平台为什么愿意自我治理？打击假货与虚假评价付出当期成本，收益却是跨期的声誉积累，平台自我治理的强度因此是一个跨期最优化问题。下述推导把这一权衡写成现值比较，给出平台自发治理的阈值条件。

\begin{derivation}{声誉外部性与自我治理}{}
设平台每期交易规模为 $t$（交易量，佣金率给定），高信任状态下每期声誉收益为 $V_H$，低信任状态下为 $V_L$，且 $V_H > V_L$：信任度高时买家愿意多交易、卖家愿意多投入，平台的每期佣金收入更高，两者的差 $V_H - V_L$ 衡量声誉对平台收入的实际价值。当期治理成本为 $C_g$（打击假货、清理虚假评价的当期投入），贴现因子为 $\delta \in (0,1)$。时机结构是关键设定：若当期治理，当期与未来各期的声誉均维持在高水平；若当期放任，当期仍可获益于存量声誉（买家不会立刻察觉），但自下期起声誉退化为低水平。这正是声誉资产的本质：可以透支，但不能无限透支。
第一步，治理的现值。选择治理时，当期承担成本但保住声誉，跨期收益为
\[
\Pi_g = (tV_H - C_g) + \delta t V_H + \delta^2 t V_H + \cdots = tV_H - C_g + \frac{\delta}{1-\delta} t V_H.
\]
第二步，放任的现值。选择放任时，当期省下治理成本、仍按高信任水平获益，但自下期起每期收益跌至 $tV_L$：
\[
\Pi_l = tV_H + \delta t V_L + \delta^2 t V_L + \cdots = tV_H + \frac{\delta}{1-\delta} t V_L.
\]
第三步，自发治理的条件。平台选择治理当且仅当 $\Pi_g \ge \Pi_l$，化简后为
\[
C_g \le \frac{\delta}{1-\delta}\, t\, (V_H - V_L).
\]
\important{当声誉的跨期价值 $\frac{\delta}{1-\delta} t (V_H - V_L)$ 不低于当期治理成本 $C_g$ 时，平台自发选择治理}。比较静态：贴现因子 $\delta$ 越高（平台越看重长期）、交易规模 $t$ 越大、信任对交易的影响 $V_H - V_L$ 越大，自我治理的激励越强；反过来，用户锁定或市场垄断使 $V_H - V_L$ 缩小，声誉约束随之失灵。
\end{derivation}

治理条件把平台的自我治理强度与其声誉资产的跨期价值绑在一起：平台的自我治理强度与声誉资产的跨期价值成正比，声誉越值钱，平台越自律。这解释了为什么成长期的平台往往比垄断期的平台更努力治理：成长期每一份声誉都对应着未来的市场份额，而垄断期的声誉流失代价被市场势力缓冲。

政策含义随之展开。监管可以利用声誉机制设计激励相容的治理制度：公开披露平台治理绩效、信用评级、白名单都是放大声誉约束的工具，让平台的治理行为进入公众与商家的决策视野，声誉的跨期价值被监管人为抬高。当声誉约束弱（市场垄断、用户锁定、经营者贴现率高）时，外部监管的必要性上升，平台治理的"自律与他律"边界由此决定：自律有效的前提是声誉值钱。中国实践提供了两者的组合形态：市场监管部门对平台企业的行政指导与约谈（第 6 章"二选一"整改）是外部约束，平台自发升级的商家治理与消费者保障机制是声誉激励的结果，两种机制在同一生态中共存互补。

\section{数字经济税收}

治理的最后一块拼图是财政层：数字经济的利润创造了价值，但既有国际税收规则在数字利润面前失能，征税权与价值创造地分离，税收竞争使各国陷入逐底竞赛。本节按"冲击、回应、方案"的顺序展开：12.4.1 节说明常设机构原则为何失效，12.4.2 节分析单边数字服务税为何只是过渡，12.4.3 节推导双支柱方案，12.4.4 节梳理实施进展与中国立场。

\subsection{常设机构原则的失效与 BEPS}

国际税收的百年规则建在一个物理前提上：利润的创造地有物理存在。规则制定的年代里，跨国经营意味着在别国设厂、开店、派驻人员，物理存在既是利润与市场国经济联系的真实载体，也是可验证的征税标志。数字交付使这个前提瓦解，本节先厘清常设机构原则及其失效的逻辑，再用转移定价算例展示税基侵蚀的具体机制。

国际税收的传统规则是：一国仅对非居民企业在本国设有\important{常设机构}（PE，Permanent Establishment）的营业利润征税，物理存在形式包括办公室、工厂、代理人。数字经济使利润的创造与物理存在脱钩：数字企业可以在无任何物理存在的情况下向市场国用户销售数字服务并获取利润，PE 规则下的征税权归零。而且这条规则被两组因素同时掏空。一是交付的虚拟化：服务以数据形态跨境交付，市场国境内的活动只剩下用户使用；二是利润的纸面移动：集团内部的转移定价可以把账面上的利润从市场国挪到低税地，后一条正是\textbf{税基侵蚀与利润转移（BEPS，Base Erosion and Profit Shifting）}问题的核心。下述算例展示转移定价的具体机制。

\begin{derivation}{BEPS 的转移定价算例}{}
设跨国集团在高税国 B 的真实营业利润为 $\pi_B = 1000$ 万元（B 国税率 $\tau_B = 30\%$），集团将无形资产（专利、品牌、算法）登记在低税国 A（$\tau_A = 5\%$），A 向 B 收取特许权使用费 $R = 800$ 万元。这种安排的商业外壳是：无形资产归 A 所有，B 使用技术当然要向 A 付费。
第一步，无转移定价安排的税负。若无形资产本就登记在 B（或使用费为零），全部利润在 B 国征税：
\[
T_0 = \tau_B \pi_B = 30\% \times 1000 = 300 \ \text{万元}.
\]
第二步，转移定价后的税负。B 国应税利润被特许权使用费扣除，降至 $\pi_B - R = 200$ 万元；A 国对使用费收入征税：
\[
T_1 = \tau_B (\pi_B - R) + \tau_A R = 30\% \times 200 + 5\% \times 800 = 60 + 40 = 100 \ \text{万元}.
\]
第三步，税基侵蚀的规模。集团总利润 1000 万元未变，总税负却从 300 万元降到 100 万元，有效税率由 $30\%$ 降至 $T_1 / \pi_B = 10\%$，\important{节税额}为
\[
T_0 - T_1 = 300 - 100 = 200 \ \text{万元}.
\]
800 万元的使用费定价只要在 A 国税率低的前提下成立，就自动完成利润搬家：合同与资产登记地的纸面安排，取代了价值创造地对利润的征税权。
\end{derivation}

算例的实质：利润跟随合同、资产与功能的纸面安排流动，而非跟随价值创造地。B 国的市场与基础设施贡献了价值，却只对剩余的 200 万元利润征税。独立交易原则（arm's length principle）要求关联交易按独立第三方的可比价格定价，这是国际税收的反避税基准；但无形资产缺乏公开市场参照价格，专利与品牌独一无二、没有现货市场，税务机关难以判定 800 万元的特许权使用费是否过高，这正是 BEPS 问题的制度根源：原则正确，执行困难。OECD 估计 BEPS 造成的税收损失约占全球企业所得税收入的 4\%--10\%（OECD，2015），被侵蚀的税基主要是高税市场国与发展中国家的税源。制度回应有两条线：单边的数字服务税（12.4.2 节）与多边的双支柱方案（12.4.3 节）。

\subsection{数字服务税的税负转嫁}

PE 失效的直接反应是市场国单边征税，产物即\textbf{数字服务税（DST，Digital Services Tax）}。各国开征 DST 的动机有三：大型数字平台的境内收入可观、多边改革进程缓慢、公众舆论对"科技巨头不缴税"高度不满。税负水平的差异是动机的量化证据：欧盟委员会估算，数字公司的平均有效税率约 9.5\%，传统企业的平均有效税率约 23\%，跨国避税安排使数字利润的系统性低税负成为各国财政的共同痛点。理解 DST 的关键不是税率本身，而是税负归宿：名义纳税人是平台，实际承担者取决于市场结构，谁最后掏钱与谁在税法上缴税是两回事。法国自 2019 年起对大型数字平台的境内数字服务收入（在线广告、用户数据销售、平台佣金）征收 3\% 的 DST，英国自 2020 年起征收 2\%，税基是收入而非利润，亏损企业同样纳税，经济性质是\textbf{流转税}而非所得税。\mn{DST 数字速记：法国 3\%（2019 年起）、英国 2\%（2020 年起）；税基三件套：在线广告、数据销售、平台佣金；性质：以收入为税基的流转税。}单边税的政治经济成本立即显现：美国认为 DST 变相针对美国企业，依据 301 条款启动调查并威胁加征报复性关税，2021 年美法达成过渡安排，法国暂缓征收，已缴税款待支柱一生效后按多边方案抵扣。单边税引发的冲突说明数字税的可持续形态只能是多边方案，而税负转嫁的分析是理解各方立场的钥匙：谁承担税负，谁就有动力推动或阻止这项税。

\begin{derivation}{DST 的税负转嫁}{}
设市场国对平台数字服务收入征收税率为 $\tau$ 的 DST，平台每单位服务向广告主收取价格 $p$，但实际只得到 $(1-\tau)p$，税收从平台的收入中直接扣除。市场均衡条件为 $D(p) = S\bigl((1-\tau)p\bigr)$，其中 $D$ 为广告主的需求，$S$ 为平台的供给。
第一步，全微分。要知道税负怎么转嫁，就考察税率从零微增一个 $d\tau$ 时均衡价格怎么变。对均衡条件两边关于 $\tau$ 全微分，并在 $\tau = 0$ 处取值（小税起始点）：
\[
D'(p)\, dp = S'\bigl[(1-\tau)\, dp - p\, d\tau\bigr]
\;\Rightarrow\;
\frac{dp}{d\tau}\Big|_{\tau=0} = \frac{S' p}{S' + |D'|}.
\]
第二步，弹性形式。用弹性改写是因为弹性无量纲、可直接比较：需求弹性 $\varepsilon_D = -D' p / Q$ 度量广告主对涨价的敏感度，供给弹性 $\varepsilon_S = S' p / Q$ 度量平台的供给反应，代入得
\[
\frac{dp}{d\tau}\Big|_{\tau=0} = \frac{p\,\varepsilon_S}{\varepsilon_S + \varepsilon_D}.
\]
第三步，税负分配。广告主支付的价格上升 $p \cdot \varepsilon_S/(\varepsilon_S + \varepsilon_D) \cdot d\tau$：\important{广告主承担的税负份额为 $\varepsilon_S/(\varepsilon_S+\varepsilon_D)$，平台自身承担的份额为 $\varepsilon_D/(\varepsilon_S+\varepsilon_D)$}。需求弹性 $\varepsilon_D$ 越小（广告主缺乏替代渠道），广告主承担的份额越大；极端情形 $\varepsilon_D \to 0$ 时 $dp = p\, d\tau$，税负全额转嫁，平台分文不损。
\end{derivation}

转嫁结果的含义：平台名义上是纳税人，但当广告主缺乏替代渠道（需求弹性小）时，税负的大部分经由广告价格转嫁给广告主，进而沿产业链传导至平台商家与消费者：\important{DST 税负的真实归宿未必是科技巨头，而可能是其产业链上的中小企业}。这个结论解释了 DST 的两大争议。其一，公平性：以收入为税基的流转税不区分盈利与亏损，亏损企业同样纳税，与所得税"有利润才缴税"的逻辑存在张力；其二，双重征税：DST 在所得税之外叠加征收，且不构成可在境外抵扣的所得税，同一笔收入被重复课税。加上单边征收引发的贸易报复，DST 的制度成本超过了它的财政收入，这两项缺陷正是 OECD 双支柱方案取代单边 DST 的理论动因。

\subsection{OECD 双支柱方案}

双支柱方案是国际税收规则百年来最大规模的重构，其必要性源于数字经济对传统税制的三重冲击：PE 原则失效（12.4.1 节）、无形资产转移定价的滥用空间、税收竞争的逐底竞赛。双支柱的分工可以概括为：支柱一回答"谁有权征税"（征税权分配），支柱二回答"最低征多少"（税率地板）。三重冲击中，最需要形式化的是第三重：税收竞争为何必然滑向逐底？

\begin{derivation}{税收竞争的逐底逻辑}{}
设两个辖区竞争同一笔流动税基 $\pi$，税率分别为 $\tau_A$、$\tau_B$，税基流向税率较低的辖区（流动成本忽略）。辖区的税收收入为
\[
T_i(\tau_i, \tau_j) = \tau_i \pi \cdot \mathbb{1}[\tau_i < \tau_j].
\]
第一步，最优反应。给定对手的税率，己方只要把税率降到对手之下一点点，就能独占整个税基：
\[
\tau_i < \tau_j \;\Rightarrow\; T_i = \tau_i \pi,\quad T_j = 0.
\]
独占的诱惑在每一轮都成立，不存在"各征各的"的稳定状态。
第二步，纳什均衡。对手的最优反应是进一步降低税率，双方反复削价，均衡趋向
\[
\tau_A = \tau_B \to 0,
\]
两辖区的税收收入趋近于零，而税基本身毫发无损：利润还在，只是没有辖区能从它身上征到税。\important{流动税基上的税率竞争是"逐底竞赛"（race to the bottom），均衡结果对所有辖区都是灾难性的}：企业坐享零税率，公共品的财源枯竭。
\end{derivation}

现实中税基并非完全流动，企业投资地的真实成本（基础设施、市场距离、法治环境）会缓冲逐底的幅度，模型刻画的是竞争的方向而非幅度；但方向本身已经说明，单边竞争没有出路，出路在于国际协调把竞争从税率维度约束到规则维度。支柱一的内容可以用一个数值算例完整呈现。

\begin{derivation}{支柱一金额 A 的数值算例}{}
设支柱一的适用门槛：全球收入超过 200 亿欧元且税前利润率超过 10\% 的跨国集团。门槛的设计意图是锁定规模最大、利润率最高的数字集团，避免把征管负担摊到普通企业头上。金额 A 的税基为剩余利润（利润率超过 10\% 的部分）的 25\%：10\% 的利润率被视为"常规回报"，留给企业所在辖区按现有规则征税；超过常规回报的剩余利润，才有 25\% 重新分配给市场国。
第一步，剩余利润。某集团全球收入 $\bar{R} = 300$ 亿欧元，税前利润率 $m = 20\%$，则剩余利润为
\[
\text{剩余利润} = (m - 10\%) \times \bar{R} = (20\% - 10\%) \times 300 = 30 \ \text{亿欧元}.
\]
第二步，金额 A 税基：
\[
\text{金额 A} = 25\% \times 30 = 7.5 \ \text{亿欧元}.
\]
第三步，按市场国收入分配。金额 A 税基按各市场国的收入占比分配：设该集团在中国市场的收入占全球收入的 15\%，中国分得的征税权为
\[
7.5 \times 15\% = 1.125 \ \text{亿欧元},
\]
中国可按本国法定税率（25\%）对该部分利润征税，可征税额约
\[
1.125 \times 25\% \approx 0.28 \ \text{亿欧元}.
\]
\important{支柱一的实质}：无论企业是否在市场国设有物理存在，市场国都按收入份额分享对剩余利润的征税权，PE 原则的缺口被制度性补上。
\end{derivation}

支柱二的计算规则同样以算例呈现。

\begin{derivation}{支柱二补足税的计算}{}
设跨国集团在某辖区的有效税率（ETR）为辖区有效税额与辖区 GloBE 所得之比：
\[
ETR_{\text{辖区}} = \frac{\text{辖区有效税额}}{\text{辖区 GloBE 所得}},
\]
全球最低税率为 $t^* = 15\%$。\important{补足税率}为最低税率与辖区有效税率之差：
\[
\text{Top-up Tax \%} = t^* - ETR_{\text{辖区}}.
\]
补足税的计税逻辑：税率是辖区的政策变量，利润落到低税辖区只缴了低税，差额部分由集团补缴，使集团在每个辖区的利润都至少按 15\% 纳税。
算例：集团在低税地的 GloBE 所得为 1000 万元，当地实缴税额 50 万元，则 $ETR = 50/1000 = 5\%$，补足税率 $= 15\% - 5\% = 10\%$，补足税额为
\[
1000 \times 10\% = 100 \ \text{万元},
\]
由母公司所在辖区（收入纳入规则 IIR）或集团其他辖区（低税支付规则 UTPR）征收。\important{支柱二为税收竞争设置地板}：低税辖区的税率低于 15\% 的部分不再产生竞争优势，利润无论落地何处都至少缴纳 15\%，逐底竞赛被制度性终止。
\end{derivation}

\begin{figure}[htbp]
\centering
\begin{tikzpicture}[>=Stealth,
    pill/.style={rectangle, rounded corners=2mm, ultra thick, align=center, font=\small, minimum width=100mm, minimum height=16mm}]
\node[pill, draw=main, fill=main!8] (P1) at (0,0) {\textbf{支柱一：征税权分配}（谁有权征）\\金额 A = 剩余利润的 25\%，按市场国收入占比分配};
\node[pill, draw=second, fill=second!8] (P2) at (0,-3.1) {\textbf{支柱二：税率地板}（最低征多少）\\补足税额 = 利润 ×（15\% − 辖区有效税率），经 IIR / UTPR 征收};
\draw[->, very thick] (P1) -- (P2);
\end{tikzpicture}
\caption{双支柱方案的分工结构}
\label{fig:pillars}
\end{figure}

图~\ref{fig:pillars} 给出了双支柱的分工结构：支柱一管征税权的分配，支柱二管税率的底线，二者针对的是不同环节的失灵：金额 A 机制针对利润征税地与价值创造地的分离，最低税机制针对辖区间的税率竞争。\mn{双支柱记忆：支柱一 = 征税权分配规则（金额 A：剩余利润的 25\% 按市场国收入分配），支柱二 = 税率底线规则（15\% 全球最低税）；补足税公式：Top-up = 利润 ×（15\% − 辖区有效税率）。}支柱一与支柱二的组合效果是：即使利润被转移至低税辖区，市场国的征税权与 15\% 的税率底线仍然共同成立，两个失灵的通道被同时封堵。

双支柱的经济影响分布在各主体之间差异显著。\textbf{低税辖区}：爱尔兰、开曼等地的"低税率招商"模式以税率差异吸引利润落地，补足税机制使利润无论落地何处都至少缴纳 15\%，低税地的竞争优势被\important{制度性取消}，其基于税收优惠的经济模式面临转型压力。\textbf{数字化企业}：大型科技集团的全球有效税率普遍低于 15\%，支柱二使其税负整体上升；支柱一则将约千亿美元量级的利润征税权从低税地重新分配给市场国（OECD 估计）。\textbf{发展中国家}：理论上受益于市场国征税权的扩张，其消费者市场首次获得对数字巨头的征税依据，但支柱一 200 亿欧元的收入门槛使绝大多数发展中国家的数字企业不在适用范围内，征管能力不足又制约其实际征收，获益取决于配套能力建设。\textbf{实施争议}：美国国内批准进程反复（支柱一触及美国科技巨头利益），欧盟以单边推进施压，全球一致性实施的时间表存在不确定性：双支柱的成败最终取决于主要经济体的政治意志而非技术细节。

\subsection{实施进展与中国立场}

规则文本的达成只是起点，双支柱的真正考验在实施。支柱二的落地进度快于支柱一，主要经济体的路径选择各不相同。补足税的执行体系由三个征收规则构成：\textbf{收入纳入规则（IIR）}由母公司所在辖区对低税子公司的利润补征，是主规则；\textbf{低税支付规则（UTPR）}在母公司辖区不实施 IIR 时由集团其他辖区兜底补征，防止母公司辖区不作为；\textbf{合格国内最低补足税（QDMTT）}由低税辖区自行按国际标准征收，使税源留在属地而非流向其他辖区，低税辖区因此有动机自征。基于实质的所得排除为真实经济活动保留空间：按有形资产与工资成本的一定比例免除补足税，避免最低税惩罚实体经营，承认厂房与雇员的真实性。进度方面，欧盟于 2022 年通过支柱二指令，2024 年起对大型集团实施 IIR，日本、韩国、英国等相继立法跟进。美国的路径特殊：未加入支柱二，以 2017 年税改引入的\textbf{全球无形资产低税收入制度（GILTI）}作为国内替代，对美国企业受控外国公司的低税利润征收最低税，有效税率由 10.5\% 升至 2026 年起的约 13.125\%；GILTI 税率低于 15\% 意味着美国企业的海外利润可能被其他辖区按 UTPR 补征，这一压力正是美国参与全球共识谈判的动力之一。

支柱一的进展慢于支柱二：金额 A 多边公约已于 2024 年开放签署，生效需足够数量的辖区完成批准；公约生效后，各国单边 DST 依约撤销，美法过渡安排也以此为终点。中国的立场清晰：支持双支柱多边方案、反对单边 DST；关注多边方案对发展中国家税源分配与征管能力的安排；中国境内数字企业大多未达支柱一门槛，但作为市场国获得的征税权增量与作为投资目的地的合规成本需要统筹评估。\important{中国的立场要点可概括为"多边、公平、可执行"：规则统一交给多边，利益分配关照发展中国家的能力约束}。这一立场与第 11 章数字贸易规则中中国"支持多边、参与区域"的一贯姿态一脉相承。

\begin{chapterbox}
\textbf{本章考点清单}

\begin{enumerate}
\item 隐私的三个经济学性质与隐私悖论（信息不对称、跨期不一致性）\important{【专硕·核心】}
\item 隐私-效用权衡的边际条件与政策比较静态 \important{【专硕·核心】}
\item 隐私规制的外部性框架与私人最优、社会最优的比较 \important{【专硕·扩展】}
\item 规制反对方四论点（信号、模式、壁垒、悖论）与 GDPR"干净数据"悖论 \important{【专硕·扩展】}
\item 三法分工（网安法/数安法/个保法）与数据分类分级 \important{【专硕·核心】}
\item 知情同意的注意力预算模型与同意疲劳 \important{【专硕·核心】}
\item 数据出境三路径与阶梯化成本 \important{【专硕·核心】}
\item 算法治理：备案、不得差异化定价、自动化决策拒绝权 \important{【专硕·扩展】}
\item GDPR 与中国制度的比较（域外管辖、被遗忘权与删除权）\important{【专硕·扩展】}
\item 避风港原则、红旗原则与电商法连带责任 \important{【专硕·核心】}
\item 平台自我治理的声誉激励条件 \important{【专硕·扩展】}
\item PE 原则失效与 BEPS 转移定价算例 \important{【专硕·核心】}
\item DST 的税负转嫁与单边税的争议 \important{【专硕·核心】}
\item 税收竞争的逐底竞赛与协调逻辑 \important{【专硕·核心】}
\item 支柱一金额 A 的算例与支柱二补足税计算 \important{【专硕·核心】}
\item 双支柱实施进展（GILTI/IIR/UTPR/QDMTT）与中国立场 \important{【专硕·扩展】}
\end{enumerate}
\end{chapterbox}

治理与税收的制度分析都建立在"识别因果关系"的实证能力之上：平台行为是否损害竞争、隐私规制是否抑制创新、数字服务税是否改变了价格结构，这些命题的检验都需要计量工具。第 13 章转入数字经济研究的方法论，为上述制度命题的实证检验提供识别策略。
